\section{Proof of Theorem~\ref{thm:gap}}\label{sec:gap}
In this section, we prove Theorem~\ref{thm:gap} with $C=2^{23}$, for $n\ge2^{512}$; that is, we assume that
\begin{equation}\label{eq:gap}
\max_\gamma |E_\gamma|\;\le\;\frac n2-2^{23}\,n^{2/3}(\log n)^{2/3}.
\end{equation}
We first prove the list edge-colouring lemma and the lemma for sharing out the edges between blocks that were described in Section~\ref{sec:sketch-gap}, and then combine them with Theorem~\ref{thm:linear}.

\subsection{A list edge-colouring lemma}\label{sec:list}
After deciding which edges between blocks are given to a block $W_i$, we must decide which tree on $W_i$ receives each of these edges.
As sketched in Section~\ref{sec:sketch-gap}, this amounts to properly edge-colouring a bipartite graph whose left vertices are the vertices outside $W_i$ and whose right vertices are colours of $\phi$, where some colours of $\phi$ may not be given to certain trees.
The following lemma does this, using Lemma~\ref{lem:BKW} twice: first for the edges at the colours with restrictions, and then for all the other edges.

\begin{lemma}\label{lem:twostage}
Let $G$ be a simple bipartite graph with vertex classes $\mathcal L$ and $\mathcal R$, such that every $v\in\mathcal L$ has degree exactly $q$ and every $\gamma\in\mathcal R$ has degree at most $q-t$. Let $\mathcal R^*\subseteq\mathcal R$, and, for each $\gamma\in\mathcal R^*$, let $J_\gamma\subseteq[q]$ be a set of size at most $L$. Suppose that every $v\in\mathcal L$ has at most $t$ neighbours in $\mathcal R^*$, and that $L\le t\le q/2$. Then, $G$ has a proper edge-colouring with colours in $[q]$ in which, for every $\gamma\in \mathcal R^*$, no edge at $\gamma$ has colour in $J_\gamma$.
\end{lemma}
\begin{proof}
First, we colour the edges at the vertices in $\mathcal R^*$, giving each such edge $v\gamma$ the list $[q]\setminus J_\gamma$, which has size at least $q-t$. In the subgraph of these edges, the vertices in $\mathcal L$ have degree at most $t\le q-t$, and the vertices in $\mathcal R^*$ have degree at most $q-t$. So, by Lemma~\ref{lem:BKW}, these edges have a proper colouring from their lists.

Next, for each $v\in \mathcal L$, let $U_v$ be the set of colours already used at $v$, so that $|U_v|\le t$. Give each remaining edge $v\gamma$ the list $[q]\setminus U_v$. In the graph of the remaining edges, each $v\in \mathcal L$ has degree $q-|U_v|$, which equals the size of its lists, and each $\gamma\in \mathcal R\setminus \mathcal R^*$ has degree at most $q-t\le q-|U_v|$. So, by Lemma~\ref{lem:BKW} again, the remaining edges have a proper colouring from their lists.

At each $v\in \mathcal L$, the two colourings use disjoint sets of colours, and each vertex in $\mathcal R$ has its edges coloured in only one of the two colourings. So, the combined colouring is proper, as required.
\end{proof}

\subsection{Distributing the edges between blocks}\label{sec:distribute}
We now show how to share out the edges between blocks, as sketched in Section~\ref{sec:sketch-gap}.
Throughout this subsection, let $V=W_1\cup\dots\cup W_k$ be a partition of the vertex set of a properly coloured complete graph, with at most one part of odd size, and let $q_i=\lfloor|W_i|/2\rfloor$ for each $i\in [k]$.
We partition each part $W_i$ of even size into $q_i$ pairs of vertices. If there is a part $W_o$ of odd size, then we choose a vertex $y\in W_o$ and partition $W_o\setminus\{y\}$ into $q_o$ pairs.
For each $i\in [k]$, we will choose a set $E_i$ of edges between $W_i$ and $V\setminus W_i$, which are the edges given to $W_i$.

We make these choices using the following $\{0,1\}$-valued variables.
For each pair $\pi=\{a,a'\}$ in $W_i$ and each pair $\pi'=\{b,b'\}$ in $W_j$ with $i<j$, the four edges between $\pi$ and $\pi'$ form a copy of $K_{2,2}$ with two perfect matchings, $P=\{ab,a'b'\}$ and $P'=\{ab',a'b\}$. We introduce a variable $z_{\pi\pi'}\in\{0,1\}$: if $z_{\pi\pi'}=1$, then we put $P$ into $E_i$ and $P'$ into $E_j$, and, if $z_{\pi\pi'}=0$, then we put $P'$ into $E_i$ and $P$ into $E_j$.
If there is a part $W_o$ of odd size, then, for each pair $\pi=\{a,a'\}$ in a part $W_i$ of even size, we introduce a variable $z_{y\pi}\in\{0,1\}$: if $z_{y\pi}=1$, then we put $ya$ into $E_i$ and do not use $ya'$, and, if $z_{y\pi}=0$, then we put $ya'$ into $E_i$ and do not use $ya$.
Let $z$ be the vector of all these variables.

The first part of the following lemma shows that the degree requirements hold for every choice of $z$. The second part shows that $z$ can be chosen so that no colour is given to any part much more often than half of the number of times it is available.

\begin{lemma}\label{lem:quads}
Let the partition $W_1,\dots,W_k$, the pairs, and the sets $E_i$ be as defined above from a vector $z$. For every choice of $z$, the following hold.
\begin{props}
\item $d_{E_i}(v)=q_i$ for every $i\in [k]$ and every $v\notin W_i$.\label{quads:deg}
\item Every edge between different parts lies in at most one set $E_i$. Such edges which lie in no set $E_i$ exist only if there is a part of odd size, in which case they are all incident with $y$, and there are exactly $q_i$ of them to each part $W_i$ of even size.\label{quads:unused}
\end{props}
Moreover, $z$ can be chosen so that, for every $i\in [k]$ and every colour $\gamma$,
\begin{equation}\label{eq:disc}
|E_i\cap E_\gamma|\le\tfrac12 c_{i,\gamma}+8,\qquad\text{where}\quad c_{i,\gamma}:=|E_\gamma\cap E(W_i,V\setminus W_i)|.
\end{equation}
\end{lemma}
\begin{proof}
For~\ref{quads:deg}, let $v\in W_j$ with $v\neq y$, and let $\pi'$ be the pair containing $v$. For each of the $q_i$ pairs $\pi$ in $W_i$, the copy of $K_{2,2}$ between $\pi$ and $\pi'$ gives exactly one edge at $v$ to $E_i$, whatever the value of $z_{\pi\pi'}$. If $v=y$ and $W_i$ has even size, then each pair in $W_i$ gives exactly one edge at $y$ to $E_i$.
For~\ref{quads:unused}, note that the copies of $K_{2,2}$ are edge-disjoint and cover all edges between different parts which are not incident with $y$, and that the variables $z_{y\pi}$ cover all such edges at $y$.

For the colour bound, let $b_{i,\gamma}(z)=|E_i\cap E_\gamma|$. This is an affine function of $z$. Its value $\beta_{i,\gamma}$ at the vector with every entry equal to $\frac12$ is half of the number of edges of colour $\gamma$ which can be put into $E_i$, so that $\beta_{i,\gamma}\le c_{i,\gamma}/2$.
A variable $z_{\pi\pi'}$ appears only in the functions $b_{i,\gamma}$ and $b_{j,\gamma}$, with coefficients $\pm(|P\cap E_\gamma|-|P'\cap E_\gamma|)$, so the sum of the absolute values of its coefficients over all the functions is at most $2\sum_\gamma(|P\cap E_\gamma|+|P'\cap E_\gamma|)=8$. For a variable $z_{y\pi}$, this sum is at most $2$.
Therefore, Lemma~\ref{lem:BF}, applied with $\Delta=8$ to the matrix of these coefficients and the vector with every entry equal to $\frac12$, gives a choice of $z$ with $|b_{i,\gamma}(z)-\beta_{i,\gamma}|\le8$ for all $i$ and $\gamma$ simultaneously, as required.
\end{proof}

We note that it is important here that a single rounding acts on all the functions $b_{i,\gamma}$ at once, as each of them depends on choices involving all the other parts. Rounding the choices between each pair of parts separately would give an error that grows with the number of parts.

We now complete the proof of Theorem~\ref{thm:gap}.
We first choose the blocks at random and record the properties we need of them (Sections~\ref{sec:blocks} and~\ref{sec:blockprops}). We then decompose the inside of each block using Theorem~\ref{thm:linear} (Section~\ref{sec:inside}), before giving each tree its edges outside its block using Lemmas~\ref{lem:quads} and~\ref{lem:twostage} (Section~\ref{sec:outside}).

\subsection{Parameters and blocks}\label{sec:blocks}
Let $n\ge2^{512}$, so that $\log n\ge354$, and let $\phi$ be a proper colouring of $K_n$ satisfying~\eqref{eq:gap}. Let $V=V(K_n)$, and let
\begin{gather*}X=(n\log n)^{1/3},\qquad B=2\lceil X/2\rceil,\qquad\varepsilon=2^{24}n^{-1/3}(\log n)^{2/3},\\ L=\lceil8\log n\rceil+2,\qquad \ell=\lceil64\log n\rceil\qquad\text{and}\qquad t=1800\ell.\end{gather*}
(The values $L$ and $t$ will be used in Lemma~\ref{lem:twostage}, and are unrelated to the functions $L(N)$ and $t(N)$ in Section~\ref{sec:recursion}.)
Then $\varepsilon X=2^{24}\log n$ and $\varepsilon\le2^{-140}$, and~\eqref{eq:gap} says precisely that $|V(\gamma)|\le(1-\varepsilon)n$ for every colour $\gamma$.

Let $k=\lfloor n/B\rfloor$. Partition $[n]$ (or $[n-1]$, if $n$ is odd) into $\lfloor n/2\rfloor$ pairs, and distribute these pairs among $k$ sets as evenly as possible. Since $kB\le n$ and $B^2\le 2(n-B)$, each set receives $B/2$ or $B/2+1$ pairs. If $n$ is odd, add the element $n$ to one of the sets. This gives a partition of $[n]$ into $k$ sets with sizes $w_i\in[B,B+3]\subseteq[X,2X]$, $i\in [k]$, at most one of which has odd size.
Let $\sigma:V\to[n]$ be a uniformly random bijection, and, for each $i\in [k]$, let the \emph{block} $W_i$ be the set of vertices which $\sigma$ maps into the $i$th of these sets. Then each $W_i$ is a uniformly random subset of $V$ of size $w_i$.
Unlike in Section~\ref{sec:recursion}, we do not condition on any event here.

\subsection{Properties of the random blocks}\label{sec:blockprops}
We show that, with positive probability, the following hold.
\begin{props}
\item Every colour has at most $L$ edges inside every block $W_i$.\label{R1}
\item $|W_i\setminus V(\gamma)|\ge\varepsilon w_i/2$ for every $i\in [k]$ and every colour $\gamma$.\label{R2}
\item For every $i\in [k]$ and every $v\notin W_i$, the number of $u\in W_i$ with $\phi(vu)\in\phi(K[W_i])$ is less than $t$.\label{R3}
\end{props}
That is, by~\ref{R1}, Theorem~\ref{thm:linear} can be applied inside each block; by~\ref{R2}, each colour misses a positive proportion of each block, which is where we use~\eqref{eq:gap}; and, by~\ref{R3}, few of the edges from a vertex to a block have a colour already appearing inside the block.

\smallskip

\noindent\textbf{Property~\ref{R1}.}
Fix $i$ and a colour $\gamma$, and let $Y$ be the number of edges of colour $\gamma$ inside $W_i$. By Lemma~\ref{lem:hyp}(a), $\E\binom Yj\le\binom{n/2}{j}(w_i/n)^{2j}\le\mu^j/j!$ with $\mu=w_i^2/(2n)\le 2X^2/n<1$. Hence $\E2^Y\le e$, so that $\Prb(Y\ge L)\le e2^{-L}\le n^{-5.5}$. A union bound over the fewer than $n^3$ pairs $(i,\gamma)$ shows that~\ref{R1} fails with probability at most $n^{-2.5}$.

\smallskip

\noindent\textbf{Property~\ref{R2}.}
Fix $i$ and a colour $\gamma$, and let $D_\gamma=V\setminus V(\gamma)$, which has size at least $\varepsilon n$. Then $\mu:=w_i|D_\gamma|/n\ge\varepsilon w_i$, and so Lemma~\ref{lem:hyp}(b) shows that $|W_i\cap D_\gamma|<\varepsilon w_i/2$ with probability at most $e^{-\varepsilon X/8}=n^{-2^{21}}$.

\smallskip

\noindent\textbf{Property~\ref{R3}.}
Fix $v$, and let $\cH_v$ be the $3$-uniform hypergraph with vertex set $V\setminus\{v\}$ whose edges are the sets $\{x,a,b\}$ for which $\phi(vx)=\phi(ab)$. As $\phi$ is proper, the edges $vx$ and $ab$ are disjoint, so these three vertices are distinct. We have $|\cH_v|\le\sum_x(|E_{\phi(vx)}|-1)<n^2/2$.
Any two vertices $p$ and $p'$ lie together in at most $3$ edges of $\cH_v$. Indeed, if $p$ plays the role of $x$, or $p'$ plays the role of $x$, or $\{p,p'\}=\{a,b\}$, then, in each case, the third vertex is determined because $\phi$ is proper.
Consequently, for each vertex $p$, the link graph $\mathrm{Lk}_v(p)$ (whose edges are the pairs $\{p',p''\}$ with $\{p,p',p''\}\in \cH_v$) has maximum degree at most $3$. It also has at most $(n/2-1)+(n-2)<3n/2$ edges, as there are fewer than $n/2$ edges of $\cH_v$ in which $p$ plays the role of $x$, and at most $n-2$ in which $p\in\{a,b\}$, since then the other vertex of $\{a,b\}$ determines the edge of $\cH_v$.

If $u\in W_i$ satisfies $\phi(vu)=\phi(ab)$ for some $a,b\in W_i$, then $\{u,a,b\}$ is an edge of $\cH_v[W_i]$, and each such edge arises in this way from at most $3$ vertices $u$. So, to show~\ref{R3}, it suffices to show that $|\cH_v[W_i]|<600\ell$ for every $i$ and every $v\notin W_i$.

Suppose first that, for some vertex $p$, the link graph $\mathrm{Lk}_v(p)$ has at least $200$ edges inside $W_i$. Since this link graph has maximum degree at most $3$, we can greedily find $40$ vertex-disjoint such edges, so that $80$ specified vertices lie in $W_i$. By Lemma~\ref{lem:hyp}(a) and a union bound over the choices of $v$, $p$ and $i$, and of $40$ edges of the link graph, this happens with probability at most
\[n^3\binom{\lfloor3n/2\rfloor}{40}\Bigl(\frac{2X}n\Bigr)^{80}\le\frac{1.5^{40}2^{80}}{40!}\,n^{3+40-80}X^{80}\le n^{-31/3}(\log n)^{80/3}\le n^{-9}.\]
Otherwise, every vertex has degree at most $199$ in $\cH_v[W_i]$. Then, if $|\cH_v[W_i]|\ge600\ell$, there are $\ell$ vertex-disjoint edges of $\cH_v$ inside $W_i$, as each chosen edge meets at most $3\cdot199<600$ edges of $\cH_v[W_i]$. Taking a union bound over $v$, $i$ and the choice of these $\ell$ edges, this happens with probability at most
\[n^2\binom{\lfloor n^2/2\rfloor}{\ell}\Bigl(\frac{2X}{n}\Bigr)^{3\ell}\le n^2\Bigl(\frac{4eX^3}{n\ell}\Bigr)^\ell=n^2\Bigl(\frac{4e\log n}{\ell}\Bigr)^\ell\le n^2(e/16)^{64\log n}\le n^{-100}.\]

The probabilities of failure for~\ref{R1}--\ref{R3} sum to less than $1$, so we can fix blocks $W_1,\ldots,W_k$ satisfying~\ref{R1}--\ref{R3}.

\subsection{Trees inside the blocks}\label{sec:inside}
By~\ref{R1}, each $K[W_i]$ has at most $L\le 8\log n+3\le2^{-49}X\le2^{-49}w_i$ edges of each colour, and $w_i\ge X\ge2^{128}$.
So, by Theorem~\ref{thm:linear}, for each block $W_i$ of even size, $E(K[W_i])$ decomposes into $q_i=w_i/2$ rainbow spanning trees $Q_{i,1},\dots,Q_{i,q_i}$ of $K[W_i]$.
If there is a block $W_o$ of odd size, then we choose $y\in W_o$ and apply the odd case of Theorem~\ref{thm:linear} with the prescribed vertex $y$. This gives $q_o$ edge-disjoint rainbow spanning trees $Q_{o,1},\ldots,Q_{o,q_o}$ of $K[W_o]$, such that the edges of $K[W_o]$ not used by these trees are the $q_o$ remaining edges at $y$ in $K[W_o]$.

\subsection{Adding the edges between the blocks}\label{sec:outside}
We apply Lemma~\ref{lem:quads} with the blocks $W_1,\ldots,W_k$, with arbitrary pairings inside the blocks, and, if there is a block of odd size, with the vertex $y$ chosen above. This gives sets $E_i$, $i\in [k]$, satisfying~\ref{quads:deg},~\ref{quads:unused} and~\eqref{eq:disc}.
By~\ref{R2}, we have $c_{i,\gamma}\le|V(\gamma)\cap W_i|\le(1-\varepsilon/2)w_i$ for each $i\in[k]$ and colour $\gamma$, so that, by~\eqref{eq:disc},
\[b_{i,\gamma}:=|E_i\cap E_\gamma|\le\frac{w_i}2-\frac{\varepsilon w_i}4+8\le q_i-t,\]
because $t+8.5\le1800(64\log n+1)+9\le2^{17}\log n\le\varepsilon X/4=2^{22}\log n$.

Now fix $i\in [k]$. Let $G_i$ be the bipartite graph with vertex classes $V\setminus W_i$ and the set of colours of $\phi$, which, for each edge $vu\in E_i$ with $v\notin W_i$, has the edge $v\,\phi(vu)$. As $\phi$ is proper, $G_i$ is simple.
By Lemma~\ref{lem:quads}\ref{quads:deg}, each vertex $v\in V\setminus W_i$ has degree $q_i$ in $G_i$, and each colour $\gamma$ has degree $b_{i,\gamma}\le q_i-t$ in $G_i$.
Let $\mathcal{R}^*=\phi(K[W_i])$ be the set of colours appearing inside $W_i$, and, for each $\gamma\in \mathcal{R}^*$, let $J_\gamma=\{j\in [q_i]:\gamma\in\phi(Q_{i,j})\}$ be the set of trees on $W_i$ containing $\gamma$. Then $|J_\gamma|\le L$ by~\ref{R1}, and each $v\in V\setminus W_i$ has fewer than $t$ neighbours in $\mathcal{R}^*$ by~\ref{R3}. Also, $L\le t\le q_i/2$, since $q_i/2\ge (X-1)/4$.
Thus, by Lemma~\ref{lem:twostage}, $G_i$ has a proper edge-colouring $c:E(G_i)\to[q_i]$ in which no edge at a colour $\gamma\in\mathcal{R}^*$ receives a colour in $J_\gamma$.
For each $j\in [q_i]$, let
\[T_{i,j}=Q_{i,j}\cup\{vu\in E_i:\ v\notin W_i\text{ and }c(v\,\phi(vu))=j\}.\]
Each $v\notin W_i$ is incident to exactly one of the edges added to $Q_{i,j}$, because $v$ has degree $q_i$ in $G_i$ and $c$ is a proper colouring with $q_i$ colours. So, $T_{i,j}$ is a spanning tree of $K_n$.
The added edges have distinct colours, since $c$ is proper at the vertices of $G_i$ which are colours of $\phi$.
Moreover, the added edges have colours which do not appear on $Q_{i,j}$: a colour in $\mathcal R^*$ avoids $J_\gamma$, and a colour not in $\mathcal{R}^*$ does not appear inside $W_i$ at all. Therefore, $T_{i,j}$ is rainbow.

\subsection{Completing the proof}\label{sec:counting}
By Lemma~\ref{lem:quads}\ref{quads:unused}, and as the trees $Q_{i,j}$ are edge-disjoint, the trees $T_{i,j}$ are pairwise edge-disjoint, and there are $\sum_iq_i=\lfloor n/2\rfloor$ of them.
If $n$ is even, every edge of $K_n$ is used. If $n$ is odd, then the unused edges are the $q_o$ edges at $y$ inside $W_o$ and, for each block $W_i$ of even size, $q_i$ edges from $y$ to $W_i$. These are $q_o+\sum_{i\ne o}q_i=(n-1)/2$ edges, all incident with $y$, and hence they form a star, as required. \qed

\begin{remark}[Where the exponent $2/3$ comes from]\label{rem:twothirds}
The proof above uses the gap in~\eqref{eq:gap} through the term $\varepsilon w_i/4$ in Section~\ref{sec:outside}. This must exceed the bound $t$ in~\ref{R3}, which is of order $X^3/n+\log n$, and the blocks must be large enough that $\varepsilon X\gtrsim\log n$, so that~\ref{R2} holds with high probability. The choice $X=(n\log n)^{1/3}$ balances these two requirements, and gives $\varepsilon\asymp n^{-1/3}(\log n)^{2/3}$. We do not know whether this exponent is a limit of our methods, or only of this particular argument.
\end{remark}
